\documentclass[10pt]{article}
\usepackage[T1]{fontenc}
\usepackage[utf8]{inputenc}
\usepackage[french]{babel}
\usepackage{amsmath,amssymb,mathrsfs,mathpazo}
\usepackage{geometry}
\geometry{a4paper,margin=2cm}
\usepackage{graphicx,array,tabularx,tabularray,colortbl,multirow,multicol,booktabs,enumitem,pifont,ulem,fancybox,diagbox,cancel}
\usepackage{pstricks,pstricks-add,pst-plot,pst-tree,pst-node,pst-text,pst-func,pst-math,pst-bspline,pst-3dplot,pst-eucl,pst-solides3d}
\let\pscalculate\undefined
\usepackage{pst-bezier}
\usepackage{tikz,pgfplots,tkz-tab}
\usetikzlibrary{arrows,arrows.meta,patterns,calc,positioning,trees,decorations.markings,shapes.geometric,shapes.misc,shapes.arrows,intersections,decorations.pathmorphing,decorations.pathreplacing,backgrounds,shadows,babel}
\pgfplotsset{compat=1.18}
\usepgfplotslibrary{fillbetween}
\usepackage[np]{numprint}
\usepackage{listings,makecell,pas-tableur,twoopt}
\let\barre\undefined
\usepackage[pstricks]{tablvar}
\definecolor{darkgreen}{RGB}{0,100,0}
\providecommand{\red}{\color{red}}
\providecommand{\blue}{\color{blue}}
\providecommand{\green}{\color{green}}
\providecommand{\black}{\color{black}}
\providecommand{\white}{\color{white}}
\providecommand{\gray}{\color{gray}}
\providecommand{\cyan}{\color{cyan}}
\providecommand{\magenta}{\color{magenta}}
\providecommand{\yellow}{\color{yellow}}
\providecommand{\violet}{\color{violet}}
\providecommand{\orange}{\color{orange}}
\providecommand{\euro}{\texteuro}
\providecommand{\EUR}{\texteuro}
\providecommand{\ieme}{\textsuperscript{e}}
\providecommand{\degres}{\ensuremath{^\circ}}
\providecommand{\degree}{\ensuremath{^\circ}}
\providecommand{\vv}[1]{\overrightarrow{#1}}
\providecommand{\Coord}[1]{#1}
\providecommand{\decosix}{}
\providecommand{\starredbullet}{\ensuremath{\bullet}}
\pagestyle{empty}
\setlength{\parindent}{0pt}
\usepackage{adjustbox}\DeclareUnicodeCharacter{2236}{:}\DeclareUnicodeCharacter{2190}{\ensuremath{\leftarrow}}
\providecommand{\R}{\mathbb{R}}
\providecommand{\N}{\mathbb{N}}
\providecommand{\Z}{\mathbb{Z}}
\providecommand{\Q}{\mathbb{Q}}
\providecommand{\D}{\mathbb{D}}
\providecommand{\C}{\mathbb{C}}
\providecommand{\e}{{\mathrm{e}}}
\providecommand{\vect}[1]{\overrightarrow{#1}}
\providecommand{\vectt}[1]{\overrightarrow{\mathrm{#1}}}
\providecommand{\vvt}[1]{\overrightarrow{\mathrm{#1}}}
\providecommand{\np}[2][]{{#2~#1}}
\providecommand{\vv}[1]{{\overrightarrow{#1}}}
\renewcommand{\d}{{\mathrm{d}}}
\providecommand{\degres}{{^\circ}}
\providecommand{\degre}{{^\circ}}
\providecommand{\blue}{\color{blue}}
\providecommand{\red}{\color{red}}
\providecommand{\enskip}{\quad}
\providecommand{\llbracket}{\lbrack\!\lbrack}
\providecommand{\rrbracket}{\rbrack\!\rbrack}
\providecommand{\euro}{{\text{€}}}
\providecommand{\no}{{\text{n°}}}
\providecommand{\interval}[2]{{[#1;#2]}}
\providecommand{\intervFO}[2]{{[#1;#2[}}
\providecommand{\dotfill}{\ldots}

\providecommand{\ds}{\displaystyle}
\providecommand{\textding}[1]{#1}
\providecommand{\Oij}{\(\mathrm{O};\vec{\imath},\vec{\jmath}\)}
\providecommand{\Oijk}{\(\mathrm{O};\vec{\imath},\vec{\jmath},\vec{k}\)}
\providecommand{\Ouv}{\(\mathrm{O};\vec{u},\vec{v}\)}
\providecommand{\intervFF}[2]{\ensuremath{[#1;#2]}}
\renewcommand{\euro}{\ifmmode\text{\texteuro}\else\texteuro\fi}
\DeclareUnicodeCharacter{00B0}{\ifmmode^{\circ}\else\textdegree\fi}
\DeclareUnicodeCharacter{20AC}{\ifmmode\text{\texteuro}\else\texteuro\fi}
% Énoncés : Éducation nationale. Transcription et corrigés : APMEP, Denis Vergès et contributeurs. Découpage et classement : CoopMaths, Sébastien Lozano, Éric Elter et contributeurs. Adaptation de la présentation : Maths Explorer.
% CC BY-SA : https://coopmaths.fr/www/about/
% https://coopmaths.fr/alea/static/bac/2025/tex/bac_2025_06_sujet2_asie_1_cor.tex
\begin{document}
\noindent\textit{APMEP / CoopMaths. CC BY-SA. Présentation adaptée par Maths Explorer.}\par\medskip
\medskip

\textbf{Partie A : Étude d'un exemple}

\begin{enumerate}
    \item D'après l'énoncé : $    P_M(T) = 0,999 \quad \text{et} \quad P_{\overline{M}}(T) = 0,005$

    \item  270\,000 personnes ont été infectées sur 750\,000 donc $ P(M) = \dfrac{270\,000}{750\,000} = 0,36$

    \item Voici l'arbre pondéré complété : 

\begin{tikzpicture}[line cap=round,line join=round,>=triangle 45,x=1.0cm,y=1.0cm]
\clip(0.68,2.16) rectangle (8,6.);
\draw [line width=1.pt] (1.,4.)-- (3.,4.6);
\draw [line width=1.pt] (1.,4.)-- (3.,3.4);
\draw [line width=1.pt] (3.7,4.75)-- (5.7,5.35);
\draw [line width=1.pt] (3.7,4.75)-- (5.7,4.35);
\draw [line width=1.pt] (3.7,3.25)-- (5.7,3.65);
\draw [line width=1.pt] (3.7,3.25)-- (5.7,2.65);
\draw (3.1,5) node[anchor=north west] {$M$};
\draw (3.1,3.6) node[anchor=north west] {$\overline M$};
\draw (5.8,5.75) node[anchor=north west] {$T$};
\draw (5.8,4.65) node[anchor=north west] {$\overline T  $};
\draw (5.8,3.) node[anchor=north west] {$\overline T  $};
\draw (5.8,4.) node[anchor=north west] {$T$};
\draw (1.7,5) node[anchor=north west] {$0,36$};
\draw (1.7,3.5) node[anchor=north west] {$0,64$};
\draw (4.2,5.7) node[anchor=north west] {$0,999$};
\draw (4.2,4.5) node[anchor=north west] {$0,001$};
\draw (4 ,2.8) node[anchor=north west] {$0,995$};
\draw (4.2,4) node[anchor=north west] {$0,005$};
\end{tikzpicture}

    \item $P(M \cap T) = P(M) \times P_M(T) = 0,36 \times 0,999 = 0,35964$

donc la probabilité que l'individu soit atteint et que le test soit positif est $0,36$ à $10^{-3}$ près.

    \item $M$ et $\overline{M}$ constituent une partition de l'univers donc d'après la loi des probabilités totales : 

$P(T) = P(M \cap T) + P(\overline{M} \cap T) = 0,36 \times 0,999 + 0,64 \times 0,005 = 0,35964 + 0,0032 = 0,36284$

donc la probabilité que l'individu ait un test positif est $0,363$ à $10^{-3}$ près.
    \item $    P_T(M) = \dfrac{P(M \cap T)}{P(T)} = \dfrac{0,35964}{0,36284} \approx 0,991$

    \item Puisque $P_T(M) \approx 0,991 > 0,95$, le test est considéré comme fiable.
\end{enumerate}

\medskip

\textbf{Partie B : Dépistage sur une population cible}

\begin{enumerate}
    \item Voici l'arbre pondéré complété : 

\begin{tikzpicture}[line cap=round,line join=round,>=triangle 45,x=1.0cm,y=1.0cm]
\clip(0.68,2.16) rectangle (8,6.);
\draw [line width=1.pt] (1.,4.)-- (3.,4.6);
\draw [line width=1.pt] (1.,4.)-- (3.,3.4);
\draw [line width=1.pt] (3.7,4.75)-- (5.7,5.35);
\draw [line width=1.pt] (3.7,4.75)-- (5.7,4.35);
\draw [line width=1.pt] (3.7,3.25)-- (5.7,3.65);
\draw [line width=1.pt] (3.7,3.25)-- (5.7,2.65);
\draw (3.1,5) node[anchor=north west] {$M$};
\draw (3.1,3.6) node[anchor=north west] {$\overline M$};
\draw (5.8,5.75) node[anchor=north west] {$T$};
\draw (5.8,4.65) node[anchor=north west] {$\overline T  $};
\draw (5.8,3.) node[anchor=north west] {$\overline T  $};
\draw (5.8,4.) node[anchor=north west] {$T$};
\draw (1.7,5) node[anchor=north west] {$p$};
\draw (1.7,3.5) node[anchor=north west] {$1-p$};
\draw (4.2,5.7) node[anchor=north west] {$0,999$};
\draw (4.2,4.5) node[anchor=north west] {$0,001$};
\draw (4 ,2.8) node[anchor=north west] {$0,995$};
\draw (4.2,4) node[anchor=north west] {$0,005$};
\end{tikzpicture}


    \item $M$ et $\overline{M}$ constituent une partition de l'univers donc d'après la loi des probabilités totales : 

$P(T) = P(M \cap T) + P(\overline{M} \cap T)  = p \times 0,999 + (1 - p) \times 0,005 = 0,999p + 0,005 - 0,005p = 0,994p + 0,005$

    \item $P_T(M) = \dfrac{P(M \cap T)}{P(T)} = \dfrac{p \times 0,999}{0,994p + 0,005}$

    \item Le test est fiable si $P_T(M) > 0,95$ :

$\dfrac{0,999p}{0,994p + 0,005} > 0,95 \iff 0,999p > 0,95(0,994p + 0,005) \iff 0,999p > 0,9443p + 0,00475 \iff  0,0547p > 0,00475 \iff  p > \dfrac{0,00475}{0,0547} \approx 0,087$ donc, le test est fiable si $p > 0,087$ (soit 8,7\%).

\end{enumerate}

\textbf{Partie C : Étude sur un échantillon}

$X$ suit une loi binomiale de paramètres $n$ et $p=0,36$.

 On cherche le plus petit $n$ tel que $P(X \geqslant 1) > 0,99$


$P(X \geq 1) > 0,99 \iff 1 - P(X = 0) > 0,99 \iff -P(X = 0) >- 0,01 \iff P(X = 0) < 0,01$

or  $P(X = 0) = (1 - 0,36)^n = 0,64^n$ donc  $P(X \geq 1) > 0,99 \iff 0,64^n < 0,01\iff  \ln(0,64^n) < \ln(0,01)$ car la fonction $\ln$ est croissante sur $]\,0\,,\,+\infty[$

$\iff  n \ln(0,64) < \ln(0,01) \iff n > \dfrac{\ln(0,01)}{\ln(0,64)} $ car $\ln(0,64)<0$ or $\dfrac{\ln(0,01)}{\ln(0,64)}\approx 10,32$

donc Il faut interroger au moins 11 individus pour que la probabilité qu'au moins l’un d’eux soit atteint dépasse 99\,\%.

\bigskip

\end{document}
